(1) 从第一面镜反射后的截面出发，一个往返的线性变换为 \(M=K(R_1)P(L)K(R_2)P(L)\)。直接矩阵乘法得 \(\det M=1\)，且 \(\operatorname{tr}M=2-4L/R_1-4L/R_2+4L^2/(R_1R_2)=4g_1g_2-2\)，所以半迹为 \(2g_1g_2-1\)。
(2) 实辛 \(2\times2\) 传输矩阵的本征值满足 \(\lambda_\pm+\lambda_\pm^{-1}=\operatorname{tr}M\)。光线扰动不指数增长的非退化稳定区要求 \(|\operatorname{tr}M|<2\)，等价于 \(-1<2g_1g_2-1<1\)，故 \(0<g_1g_2<1\)。等号位置有退化或边缘稳定情形，必须单独分析，不能无条件作为严格稳定腔。
(3) 对 \(R_1=R_2=2L\)，有 \(g_1=g_2=1/2\)、\(g_1g_2=1/4\)，所以 \(\operatorname{tr}M=-1\)。本征值为 \(e^{\pm i\psi}\)，\(\cos\psi=\operatorname{tr}M/2=-1/2\)，可取 \(\psi=2\pi/3\)，位于严格稳定区。
(4) 对称腔中束腰到任一镜的距离 \(z=L/2\)，基模波前曲率满足 \(R=z+z_R^2/z\)。当 \(R=2L\) 时 \(z_R^2=(L/2)(2L-L/2)=3L^2/4\)，故 \(z_R=\sqrt3L/2\)；由 \(z_R=\pi w_0^2/\lambda\) 得 \(w_0^2=\lambda\sqrt3L/(2\pi)\)。该关系要求衍射近轴、镜口径足够大；在几何稳定域边界处需额外考虑衍射损耗和退化模。
